\section{Supplementary proofs and technical analysis}\label{app:proofs}

\subsection{Conformal rank is super-uniform under ties}

\begin{lemma}[Conservative conformal rank]\label{lem:rank}
Fix action \(a\), and let \(H_1,\ldots,H_m,X\) be exchangeable values of \(S_a\). Define
\[
 p_a(X)=\frac{1+\sum_{i=1}^{m}\ind\{S_a(H_i)\geq S_a(X)\}}{m+1}.
\]
Then for every real \(u\),
\[
 \Prob\{p_a(X)\leq u\}\leq u.
\]
In particular, with \(\alpha_a=\alpha w_a\) and nonnegative weights summing to at most one, \(\sum_a\Prob\{p_a\leq\alpha_a\}\leq\alpha\).
\end{lemma}
\begin{proof}
Exchangeability implies the joint distribution of \((S_a(H_1),\ldots,S_a(H_m),S_a(X))\) is invariant under permutations. For any fixed real \(z\), the rank of \(z\) among the \(m+1\) exchangeable values,
\[
 R(z)=1+\sum_{i=0}^{m}\ind\{V_i>z\},\qquad
 (V_0,\ldots,V_m)=(S_a(X),S_a(H_1),\ldots,S_a(H_m)),
\]
is a uniformly distributed element of \(\{1,\ldots,m+1\}\) whenever the \(V_i\) are almost surely distinct. With ties, the weak upper rank
\[
 \widetilde R(z)=1+\sum_{i=0}^{m}\ind\{V_i\geq z\}-\ind\{V_0\geq z\}
\]
dominates the position in the sorted array, so
\[
 \Prob\{\widetilde R(S_a(X))\leq k\}\leq \frac{k}{m+1}.
\]
Setting \(k=\lfloor u(m+1)\rfloor\) yields the claim. The weight statement is the union bound with \(\sum_a w_a\leq1\).
\end{proof}

Weak inequalities are required for ties: using strict inequalities can concentrate ranks at one under constant scores.

\subsection{Full statement of Proposition~\ref{prop:spending}}\label{app:registered}

Let \(\calD_{\rm dev}\) denote all development choices: score maps, detectors, budgets, execution order (if any), and weights. Weights may be instantiated by a development-frozen rule depending only on \((m,\alpha)\), the registered family \(\calA\), and a development-fixed priority order---never on the calibration ranks or the test document. Conditional on \(\calD_{\rm dev}\), Assumption~\ref{ass:exchange} gives document-level exchangeability of \(H_1,\ldots,H_m,X\sim P_0\).

\begin{proposition}[Conditional on development]\label{prop:spending-full}
Under Assumption~\ref{ass:exchange} and the development freezing above,
\[
 \Prob_{P_0}\{\text{alert at any executed action}\mid \calD_{\rm dev}\}
 \leq \sum_{a\in\calA}\alpha w_a \leq \alpha.
\]
The bound holds for arbitrary data-dependent execution subsets and orders drawn from \(\calA\), and for arbitrary stopping rules measurable with respect to the ranks seen so far.
\end{proposition}
\begin{proof}
By Lemma~\ref{lem:rank}, for each fixed \(a\), \(\Prob_0(p_a\leq\alpha w_a\mid\calD_{\rm dev})\leq\alpha w_a\). The policy alerts only if some executed \(a\) satisfies \(p_a\leq\alpha w_a\). Since the executed set is a subset of \(\calA\), the alert event is contained in \(\bigcup_{a\in\calA}\{p_a\leq\alpha w_a\}\). The union bound and \(\sum_a w_a\leq1\) finish the proof. Note that the individual events \(\{p_a\leq\alpha w_a\}\) may be arbitrarily dependent; no conditional super-uniformity along the path is used.
\end{proof}

Zero-weight actions cannot alert: their threshold is zero and ranks are at least \(1/(m+1)>0\). They may still execute, for logging or routing features fixed in development.

\subsection{Complete-path consistency}\label{app:complete-path}

Proposition~\ref{prop:path} requires that deployment follow a \emph{prefix} of the frozen policy path. Formally: for each \(t\leq T_\pi(X)\),
\[
 M_{\pi,t}(X)=\max_{s\leq t} g_{a_s(X)}(S_{a_s(X)}(X))
 \leq \max_{s\leq T_\pi(X)} g_{a_s(X)}(S_{a_s(X)}(X))=M_\pi(X),
\]
so \(p^{\rm path}_t(X)\geq p^{\rm path}_{T_\pi(X)}(X)\) with the complete calibration reference. Early termination therefore cannot produce a smaller \(p\)-value than the terminal one.

\paragraph*{Empty and failed routes.}
If the route is empty, \(M_\pi=-\infty\) and no alert is issued. If every action fails, the partial and complete maxima are both \(-\infty\); the rank is then
\[
 p^{\rm path}=\frac{1+\sum_i\ind\{M_\pi(H_i)\geq -\infty\}}{m+1}=1,
\]
provided every calibration complete-path maximum is a real number (at least one successful action on each calibration path). If some calibration paths are also all-failure, the count \(\ind\{-\infty\geq-\infty\}=1\) still yields \(p=1\) when the test path is all-failure. Alerts therefore never fire on empty or fully failed routes.

\paragraph*{Induced sequential tests.}
Koning and van Meer show how to induce anytime-valid sequential tests from terminal tests under suitable embedding \citep{koning2026anytime}. Construction B is the special case where the terminal test is a conformal rank of the complete-path maximum \(M_\pi\) and the sequential monitor compares partial maxima against the same terminal reference. The containment \(M_{\pi,t}\leq M_\pi\) is exactly the monotonicity condition that makes the induced sequential \(p\)-value no smaller than the terminal one. We do not claim novelty of the induction; we specify how the terminal statistic, the calibration sample, and the policy freezing interact in this application.

\subsection{Calibration resolution for equal weights}\label{app:resolution}

With \(n\) equal-weight active actions, \(w=1/n\) and \(\alpha_a=\alpha/n\). Feasibility \eqref{eq:resolution} requires
\[
 m\geq \left\lceil\frac{n}{\alpha}\right\rceil-1.
\]
For the prototype family of sixteen actions at equal weight:
\begin{itemize}
\item \(\alpha=0.05\): \(m\geq 319\);
\item \(\alpha=0.01\): \(m\geq 1{,}599\);
\item \(\alpha=0.001\): \(m\geq 15{,}999\).
\end{itemize}
A development-frozen rule may concentrate the budget on \(k=\lfloor\alpha(m+1)\rfloor\) active actions with weight \(1/k\) each and zero elsewhere; then the active weights satisfy \(k\cdot(1/k)\leq1\) and each active action meets \eqref{eq:resolution} at the current \(m\). Remaining registered actions receive \(w=0\) and cannot alert.

\subsection{Marginal validity under a shifted test law}\label{app:shift}

Let \(\calH\sim P_{\rm cal}^{m}\) be the calibration mechanism and \(X\sim P_0\) the null test law under the design. The marginal guarantee is
\[
 \Prob_{\calH,X}\{\text{alert}\}\leq\alpha.
\]
If the test law shifts to \(\widetilde P_0\) while the calibration mechanism is unchanged, total variation over \(X\) alone gives
\[
 \Prob_{\calH,\widetilde X}\{\text{alert}\}
 \leq \Prob_{\calH,X}\{\text{alert}\}+\TV(P_0,\widetilde P_0)
 \leq \alpha+\TV(P_0,\widetilde P_0).
\]
A fixed calibration sample \(\calH=h\) can have conditional FPR
\(\Prob_X\{\text{alert}\mid \calH=h\}\) well above \(\alpha\) on a set of \(h\)'s of small but positive \(P_{\rm cal}\)-measure; marginal control averages over that measure. Deployment-time conditional claims therefore require calibration-conditional methods \citep{bates2023} or larger calibration sizes.
